% =============================================================================
% Assignments/CS301/10-file-organization-hashing-assignment.tex
% Assignment 10 (Week 10): File Organization & Hashing
%
% Student set — NO answers. Solution key:
% 10-file-organization-hashing-solutions.tex (gitignored via
% **/*-solutions.tex, never deployed). Sourced from
% Slides/CS301/10-file-organization-hashing.tex (54 pp) and the Week 10 topic
% list in syllabi/CS301.md, objective "Select a file organization and design
% a hash function with an appropriate collision-resolution scheme".
%
% Fresh instances throughout — no problem re-asks a deck example. Deck
% instances avoided: keys 12,44,13,88,23,94,11,39,20,16,5 with m=11; the
% 1,000-record 64 B file and roll 20230457; "cat" with m=97, B=31; linear
% probe trace 12,44,13,88,23,94; double hashing for k=23; multiplication
% m=16, A=0.618, k=123456; mid-square 5678; folding 123456789; digit
% analysis on the 2023 prefix.
%
% This set uses: Q3 division m=13 with keys 47,82,15,63,91,28,74,36,59 and
% multiplication m=16, A=0.618034 with k=2947,7384,8417; Q4 linear probing
% m=13 inserting 24,45,62,11,76,38,90,17,53,29, deleting 38, then inserting
% 79; Q5 a 150,000-account print-quota service in a 4 MiB slot budget;
% Q6 a 900,000-record utility meter file with a 70/30 point/range workload.
%
% Compile with (macOS/Linux): cd Assignments/CS301 &&
%   TEXINPUTS=../../Preambles:$TEXINPUTS xelatex -interaction=nonstopmode 10-file-organization-hashing-assignment.tex
%   BIBINPUTS=../..:$BIBINPUTS bibtex 10-file-organization-hashing-assignment
%   TEXINPUTS=../../Preambles:$TEXINPUTS xelatex -interaction=nonstopmode 10-file-organization-hashing-assignment.tex (x2)
% =============================================================================

\documentclass[11pt]{article}
\usepackage[margin=1in]{geometry}
% =============================================================================
% Preambles/header.tex — shared LaTeX/Beamer preamble for this project
%
% Use from a Beamer lecture by prepending:
%     \input{header}
% and compiling with TEXINPUTS=../Preambles:$TEXINPUTS (the /compile-latex
% skill does this automatically).
%
% PALETTE CONTRACT. Color names defined here MUST match the names in
% Quarto/theme-template.scss so Beamer and Quarto renderings use the same
% palette. The scripts/check-palette-sync.sh script enforces this.
%
% To customize: edit the HEX values below AND the matching $variable in
% Quarto/theme-template.scss, then run scripts/check-palette-sync.sh.
% =============================================================================

% -----------------------------------------------------------------------------
% Engine detection + encoding
%
% Our /compile-latex skill uses XeLaTeX, where inputenc is unnecessary and
% can produce encoding warnings. Only load inputenc under pdfTeX.
% -----------------------------------------------------------------------------
\usepackage{iftex}
\ifPDFTeX
  \usepackage[utf8]{inputenc}
\fi

\usepackage{amsmath, amssymb, amsthm}
\usepackage{booktabs}
\usepackage{graphicx}

% -----------------------------------------------------------------------------
% PALETTE — mirrors Quarto/theme-template.scss variable names.
% Load xcolor and define colors BEFORE anything that references them
% (notably hyperref, hypersetup, and the Beamer color assignments).
% -----------------------------------------------------------------------------
\usepackage{xcolor}

\definecolor{primary-blue}{HTML}{012169}      % headings, accents
\definecolor{primary-gold}{HTML}{B9975B}      % emphasis, borders
\definecolor{highlight-yellow}{HTML}{F2A900}  % markers, alerts
\definecolor{light-bg}{HTML}{E8EDF5}          % title / section backgrounds
\definecolor{jet}{HTML}{1A1A1A}               % body text

% Semantic colors (used in diagrams and annotations)
\definecolor{positive}{HTML}{15803D}          % good / observed / identified
\definecolor{negative}{HTML}{B91C1C}          % bad / problematic / bias
\definecolor{neutral}{HTML}{525252}           % reference / context

% Highlight accents (match .hi-* classes in the SCSS)
\definecolor{hi-slate}{HTML}{314F4F}
\definecolor{hi-green}{HTML}{2E8B57}
\definecolor{hi-red}{HTML}{C41E3A}

% Semantic aliases so skill templates and snippets can write
% \textcolor{accent}{...} without hard-coding "primary-blue".
\colorlet{accent}{primary-blue}
\colorlet{accent2}{primary-gold}

% -----------------------------------------------------------------------------
% Hyperref — loaded AFTER xcolor + palette so link colors resolve.
% -----------------------------------------------------------------------------
\usepackage{hyperref}
\hypersetup{colorlinks=true, linkcolor=primary-blue, urlcolor=primary-blue,
            citecolor=primary-blue, pdfborder={0 0 0}}

% -----------------------------------------------------------------------------
% Course code — set once per deck (after \input{header}, before \begin{document})
% via \coursecode{CS401}. Shown in the footer on every slide so a deck's course
% is identifiable without opening the title page. Empty by default.
% -----------------------------------------------------------------------------
\makeatletter
\newcommand{\coursecode}[1]{\gdef\@coursecodetext{#1}}
\gdef\@coursecodetext{}
\makeatother

% -----------------------------------------------------------------------------
% Beamer theme assignments (only apply under Beamer).
%
% We detect Beamer via \@ifclassloaded{beamer}, which is the documented,
% non-internal way. This must be wrapped in \makeatletter/\makeatother
% because \@ifclassloaded contains an @-catcode character.
% -----------------------------------------------------------------------------
\makeatletter
\@ifclassloaded{beamer}{%
  \usefonttheme{professionalfonts}
  \setbeamercolor{structure}{fg=primary-blue}
  \setbeamercolor{frametitle}{fg=primary-blue}
  \setbeamercolor{title}{fg=primary-blue}
  \setbeamercolor{subtitle}{fg=primary-gold}
  \setbeamercolor{author}{fg=primary-blue}
  \setbeamercolor{institute}{fg=neutral}
  \setbeamercolor{date}{fg=primary-gold}
  \setbeamercolor{itemize item}{fg=primary-blue}
  \setbeamercolor{itemize subitem}{fg=primary-gold}
  \setbeamercolor{alerted text}{fg=primary-gold}
  \setbeamercolor{block title}{fg=white, bg=primary-blue}
  \setbeamercolor{block body}{bg=light-bg}
  % Semantic block variants: block=definition/concept (blue), exampleblock=
  % worked example (green/positive), alertblock=key takeaway or caution (gold).
  % Keep to <=2 colored boxes per slide across ALL three variants (INV-7).
  \setbeamercolor{block title example}{fg=white, bg=positive}
  \setbeamercolor{block body example}{bg=light-bg}
  \setbeamercolor{block title alerted}{fg=white, bg=primary-gold}
  \setbeamercolor{block body alerted}{bg=light-bg}

  % Minimal footer: course code (left) + page X / N (right), no navigation clutter
  \setbeamertemplate{navigation symbols}{}
  \setbeamertemplate{footline}{%
    \color{neutral}\footnotesize\hspace{0.7em}\@coursecodetext\hfill
    \insertframenumber/\inserttotalframenumber\hspace{0.7em}\vspace{0.3em}}
}{}
\makeatother

% -----------------------------------------------------------------------------
% TikZ setup — libraries the snippet gallery uses
% -----------------------------------------------------------------------------
\usepackage{tikz}
\usetikzlibrary{arrows.meta, positioning, calc, decorations.pathreplacing,
                fit, shapes.geometric, backgrounds}

% Shared TikZ styles — snippets and hand-written diagrams can pick these up
% via `\begin{tikzpicture}[every node/.style={...}]` or by naming specific
% styles (e.g., `\node[dag-node] (X) at (0,0) {X};`).
\tikzset{
  % Node styles for causal DAGs and similar
  dag-node/.style = {
    draw=primary-blue, fill=light-bg, rounded corners=2pt,
    minimum width=1.2cm, minimum height=0.9cm, align=center, font=\small
  },
  decision-node/.style = {
    draw=primary-gold, fill=white, diamond, aspect=1.8,
    minimum width=1.8cm, minimum height=1.2cm, align=center, font=\small,
    inner sep=1pt
  },
  % Edge styles — observed vs counterfactual
  observed-edge/.style = {-{Latex[length=2.5mm]}, thick, draw=jet},
  counterfactual-edge/.style = {-{Latex[length=2.5mm]}, thick, draw=neutral, dashed},
  confound-edge/.style = {-{Latex[length=2.5mm]}, thick, draw=negative, dashed},
  % Data-point styles for plots
  observed-dot/.style = {fill=primary-blue, draw=primary-blue, circle, inner sep=1.2pt},
  counterfactual-dot/.style = {fill=white, draw=neutral, circle, inner sep=1.2pt}
}

% -----------------------------------------------------------------------------
% Convenience macros
% -----------------------------------------------------------------------------
\newcommand{\muted}[1]{\textcolor{neutral}{#1}}
\newcommand{\key}[1]{\textcolor{primary-gold}{\textbf{#1}}}
\newcommand{\good}[1]{\textcolor{positive}{#1}}
\newcommand{\bad}[1]{\textcolor{negative}{#1}}

% Standout frame macro: full-bleed dark background for section transitions.
% Beamer-only (uses \begin{frame}/\setbeamercolor/\usebeamerfont) -- do not
% call from an article-class file (e.g. Notes/<CODE>/*-notes.tex). \muted,
% \key, \good, \bad, and \coursecode above are plain xcolor/gdef and are
% safe to use from either class; verified when Notes/article-class support
% was added.
% Expand: \transitionslide{Your transition title}
\newcommand{\transitionslide}[1]{%
  \begingroup
  \setbeamercolor{background canvas}{bg=primary-blue}
  \begin{frame}[plain]
    \centering
    \vfill
    {\usebeamerfont{title}\color{white}\Large #1\par}
    \vfill
  \end{frame}
  \endgroup
}

% Section divider frame (Beamer-only), an alternative to \transitionslide with a
% darker two-line style: near-black (jet) background, a small white label line
% above a larger gold title, and a thin gold rule along the bottom edge.
% Expand: \sectiondivider{Section 2}{Pointers}
\newcommand{\sectiondivider}[2]{%
  \begingroup
  \setbeamercolor{background canvas}{bg=jet}
  \begin{frame}[plain]
    \vspace*{3.0cm}
    {\color{white}\ttfamily\large #1\par}
    \vspace{0.5cm}
    {\color{primary-gold}\LARGE #2\par}
    \begin{tikzpicture}[remember picture, overlay]
      \draw[primary-gold, line width=2.5pt]
        ($(current page.south west)+(0,0.1)$) -- ($(current page.south east)+(0,0.1)$);
    \end{tikzpicture}
  \end{frame}
  \endgroup
}

% -----------------------------------------------------------------------------
% Channel watermark — "Concepts That Click" (YouTube).
%
% Article-class files (CompetitiveExam Q/A sets, Notes, Assignments) get a
% light diagonal draft watermark on every page plus a centered footer naming
% the channel. Beamer decks get a subtle TikZ background watermark instead
% (the footline already carries the course code). Centralized here so every
% MCQ PDF is branded without per-file edits.
% -----------------------------------------------------------------------------
\makeatletter
\@ifclassloaded{article}{%
  \usepackage{draftwatermark}
  \SetWatermarkText{Concepts That Click}
  \SetWatermarkScale{1}
  \SetWatermarkFontSize{2cm}
  \SetWatermarkLightness{0.86}
  \SetWatermarkAngle{45}
  \usepackage{fancyhdr}
  \pagestyle{fancy}
  \fancyhf{}
  \fancyfoot[C]{\footnotesize\textcolor{neutral}{Concepts That Click~|~YouTube: \href{https://www.youtube.com/@concepts.thatclick}{@concepts.thatclick}}}
  \renewcommand{\headrulewidth}{0pt}
  \renewcommand{\footrulewidth}{0pt}
  \fancypagestyle{plain}{%
    \fancyhf{}%
    \fancyfoot[C]{\footnotesize\textcolor{neutral}{Concepts That Click~|~YouTube: \href{https://www.youtube.com/@concepts.thatclick}{@concepts.thatclick}}}%
    \renewcommand{\headrulewidth}{0pt}%
    \renewcommand{\footrulewidth}{0pt}%
  }
}{}
\@ifclassloaded{beamer}{%
  \setbeamertemplate{background}{%
    \begin{tikzpicture}[remember picture, overlay]
      \node[rotate=45, opacity=0.055, align=center]
        at (current page.center)
        {\Huge\textcolor{neutral}{Concepts That Click}};
    \end{tikzpicture}%
  }
}{}
\makeatother 
\coursecode{CS301}
\renewcommand{\thesection}{10.\arabic{section}}

% Emergency stretch so long \texttt{} tokens in prose do not spill the margin.
\setlength{\emergencystretch}{2.5em}

\title{Assignment 10 (Week 10): File Organization \& Hashing}
\author{Auqib Hamid Lone\\[0.3em]
  \emph{Department of Computer Science \& Engineering}, \emph{GCET Ganderbal}}
\date{\today}

\begin{document}
\maketitle

\noindent\textbf{Due:} two weeks from issue. There is no automatic late penalty --- if you need more time, ask before the deadline and an extension will be granted (see the course policy in the syllabus).

\noindent Answer every question. Show all working for Numerical and Design questions. Each problem states its own assumptions; everything you need is in Week~10's lecture alone. Notation matches the lecture exactly: a \key{record} is a fixed-length collection of fields and its \key{key} is the field a search compares; a \key{file} is a collection of records on a device; the cost of a file operation is counted in \emph{block transfers}, not key comparisons. A hash table stores $n$ keys in $m$ slots through a \key{hash function} $h(k)$, and the \key{load factor} is $\alpha=n/m$. A \key{probe} is one slot examined. \key{Open hashing} means \key{separate chaining} (colliding keys live in lists outside the table); \key{closed hashing} means \key{open addressing} (every key lives in one of the $m$ slots). For reference, the lecture's expected-probe formulas under uniform hashing are
\[
  \text{chaining: } 1+\frac{\alpha}{2}\ \text{(hit)},\quad \alpha\ \text{(miss)};
  \qquad
  \text{linear probing: } \frac{1}{2}\Bigl(1+\frac{1}{1-\alpha}\Bigr)\ \text{(hit)},\quad
  \frac{1}{2}\Bigl(1+\frac{1}{(1-\alpha)^{2}}\Bigr)\ \text{(miss)}.
\]
Whenever a problem asks for a structure choice, name the workload property that decides it and give the cost that the property turns on.

\section{Conceptual}

\textbf{Q1.} A file organization is a decision about where records live; an access scheme is a decision about how the device may be read (Horowitz \& Sahni Ch.~8; Aho--Hopcroft--Ullman Ch.~4) \cite{HorowitzSahni2008_fundamentals_data_structures,AhoHopcroftUllman1983_data_structures_algorithms}.
\begin{enumerate}
  \item[(a)] Name the four access schemes from the lecture. For each, state how the device is used, the medium it assumes, and whether it can jump straight to an arbitrary record. Then explain why a file operation is priced in block transfers rather than in key comparisons, and why the same search algorithm can have a different cost on tape and on disk.
  \item[(b)] The lecture names four file organizations: sequential (in both \emph{pile} and \emph{sorted} variants), direct (relative), indexed sequential, and hashed. For each, give the search, insert, and range-scan costs from the Week~10 comparison table, and name one workload for which that organization is the right choice. State which organization is the only one that is strong in both point lookups and range scans, and why.
  \item[(c)] A district hospital keeps $250{,}000$ patient records of $128$ bytes each, keyed by a $10$-digit national health ID. Only $250{,}000$ of the $10^{10}$ possible ID values are in use, and they are scattered. The disk transfers one $4$\,KiB block per seek, and the records pack $32$ to a block. The workload is: insert on admission, delete on discharge, and a point lookup by ID at the triage desk. There are no range queries, no ordered reports, and no batch pass over all records. Choose one of the four organizations and justify the choice against each of the other three using these facts --- not with the general claim that it is ``faster''. Give the block count of the file and the cost in block transfers that decides the case.
\end{enumerate}
Why this matters: every later database decision is this same argument in a larger setting --- the organization is chosen by the workload's access pattern, and the price is paid in device blocks, so naming the pattern correctly is the whole skill.

\textbf{Q2.} The two families of collision resolution carry names that mislead, and the price of hashing is paid in the operations it drops (Sedgewick Sec.~3.4, p.~458; Wirth Sec.~5.1, p.~177) \cite{Sedgewick2011_algorithms,Wirth2004_algorithms_data_structures}.
\begin{enumerate}
  \item[(a)] Define \key{open hashing} and \key{closed hashing}, giving the alternative name for each. State which family keeps colliding keys \emph{outside} the $m$ slots, which keeps every key \emph{inside} them, which permits $\alpha>1$, and which needs a tombstone to delete. Then explain what ``open'' and ``closed'' originally referred to, and why an unqualified use of either name is a mistake in an exam answer.
  \item[(b)] Hashing buys an expected $O(1)$ point lookup and gives up order. Name the four ordered operations the table can no longer answer directly, and explain the structural reason: which properties of $h$ and of the slot layout are missing? Then price the workaround --- for a hashed file of $n$ keys in $m$ buckets, what does ``report every key with value in $[a,b]$'' cost, and why can that scan not stop early once it has passed $b$?
  \item[(c)] A student writes: ``Closed hashing means the table is closed, so once it is full no further keys can be inserted; and since a lookup is $O(1)$, a range query is $O(1)$ too once the two endpoints are known.'' Identify the two errors and correct each in one or two sentences.
\end{enumerate}
Why this matters: the names are examined every year, and the lost order is the reason Week~11 exists --- if you cannot say precisely what hashing gives up, you cannot justify choosing it.

\section{Numerical}

\textbf{Q3.} Hash arithmetic and the price of a chain (Wirth Sec.~5.1, p.~177; Sec.~5.3, p.~178; Sedgewick Sec.~3.4, p.~458) \cite{Wirth2004_algorithms_data_structures,Sedgewick2011_algorithms}.
\begin{enumerate}
  \item[(a)] Use the division method with $m=13$, so $h(k)=k \bmod 13$, on the nine keys
  \[
    47,\ 82,\ 15,\ 63,\ 91,\ 28,\ 74,\ 36,\ 59,
  \]
  inserted in that order into a chaining table with one list per slot and insertion at the \emph{tail} (so each chain stays in insertion order). Give $h(k)$ for every key, draw the final table, give every chain length, and state how many of the nine keys collided. Then compute the \emph{exact} average number of key comparisons for a successful search over these nine keys, and compare it with the uniform-hashing estimate $1+\alpha/2$ at $\alpha=9/13$. Explain the direction of the difference.
  \item[(b)] Use the multiplication method with $m=16$ and $A=0.618034$:
  \[
    h(k)=\bigl\lfloor 16\,(kA \bmod 1)\bigr\rfloor .
  \]
  Compute $h(k)$ for $k=2947$, $k=7384$, and $k=8417$, showing $kA$ to three decimal places and the fractional part you use. Explain why $m=16$ is acceptable for this method but would be a poor choice for the division method, and say which bits of $k$ decide the slot under each of the two methods.
  \item[(c)] The same index now holds $n=20$ keys. Using the lecture's chaining formulas, give $\alpha$, the expected probes for a successful search, and the expected probes for an unsuccessful search when $m=13$. Recompute all three for $m=41$. Which table would you ship, and what does this pair of calculations show about the role of $m$ when $n$ is fixed?
\end{enumerate}
Why this matters: the division and multiplication methods are the two the exam asks you to evaluate by hand, and the chaining formula is the only place where the load factor becomes an actual number of probes rather than a slogan.

\textbf{Q4.} A linear-probing trace, including the deletion that breaks it (Sedgewick Sec.~3.4, p.~458; Wirth Sec.~5.4, p.~180) \cite{Sedgewick2011_algorithms,Wirth2004_algorithms_data_structures}.
A table has $m=13$ slots, $h'(k)=k \bmod 13$, and resolves collisions by linear probing, $h(k,i)=(h'(k)+i) \bmod 13$ for $i=0,1,2,\dots$. Insert the keys
\[
  24,\ 45,\ 62,\ 11,\ 76,\ 38,\ 90,\ 17,\ 53,\ 29
\]
in that order, starting from an empty table.
\begin{enumerate}
  \item[(a)] For each key, give its home slot $h'(k)$, the full probe sequence, the number of probes, and the slot it finally occupies. Draw the final $13$-slot table, state the total number of probes spent on the ten inserts, and give the final load factor $\alpha$.
  \item[(b)] Delete $38$ and mark its slot \texttt{DELETED}. Give the probe sequence and the probe count for a search for $90$ after the deletion. Then state what a search for $90$ would have concluded if slot $1$ had instead been cleared to \texttt{EMPTY}, and why it would have been wrong.
  \item[(c)] Before inserting anything else, give the probe sequence and the probe count for an unsuccessful search for $50$, and explain why the search cannot stop at the tombstone. Then insert $79$ using the lecture's rule that an insert may reuse the first tombstone it meets: give the home slot, the probes examined, and the slot it occupies. Finally, state where $79$ would have gone if tombstones could not be reused.
\end{enumerate}
Why this matters: deletion is the one operation where open addressing has a correctness rule rather than a performance rule --- a table that answers ``absent'' for a key it is holding is a bug no load factor can excuse.

\section{Design}

\textbf{Q5.} Design the hash index for a campus print-quota service (Sedgewick Sec.~3.4, p.~458; Karumanchi Ch.~14, pp.~610--639 PDF) \cite{Sedgewick2011_algorithms,Karumanchi2017_data_structures_algorithms_made_easy}. The service holds $150{,}000$ active accounts, keyed by a $9$-character roll string such as \texttt{2023CSE01}. Lookups run on every print job, peaking near $8{,}000$ per second; accounts are inserted at registration and deleted in a burst at the end of each semester. The index is an array of $m$ slots indexed $0,\dots,m-1$, each slot costs $16$ bytes, and the service may spend at most $4$\,MiB on that array (any chaining nodes would be extra, outside this budget). The lookup endpoint is public; the roll strings themselves are assigned by the registry.
\begin{enumerate}
  \item[(a)] \emph{Sizing.} From the $4$\,MiB slot budget and the $16$-byte slot, compute the largest possible $m$. Choose $m$ as the largest prime not exceeding that bound, and state it. Give $\alpha$ at $150{,}000$ accounts, and the expected probes for a hit and for a miss under (i) chaining and (ii) linear probing at that $\alpha$. Which policy do the probe counts favour, and what does the memory budget say about it?
  \item[(b)] \emph{The function.} Specify $h$ for a $9$-character string as a polynomial in base $B=31$, reduced modulo $m$ at every step. Show what goes wrong if $m$ were $2^{p}$ instead of a prime. Then decide whether this fixed $h$ is safe for this service, given that the endpoint is public but the keys are registry-assigned; state the condition under which you would replace it with a universal family, and give that family's form.
  \item[(c)] \emph{Policy.} State the collision policy you would ship, the load-factor threshold that triggers a resize, the resize target, and the deletion and rehash plan. Justify the threshold with the probe counts from part (a). Finally, state what the first resize does to the memory budget, and when in the semester's workload you would schedule it.
\end{enumerate}
Why this matters: an $O(1)$ promise is a promise about $\alpha$ and about the key distribution --- this is the question that makes you turn ``use a hash table'' into numbers a colleague could implement and an operator could run.

\textbf{Q6.} Design the file organization for a power utility's meter file (Horowitz \& Sahni Ch.~8; Wirth Sec.~5.1, p.~177) \cite{HorowitzSahni2008_fundamentals_data_structures,Wirth2004_algorithms_data_structures}. The file holds $900{,}000$ premises records of $96$ bytes each, keyed by premises ID. IDs are integers assigned in increasing blocks by pincode and street, but the values are sparse: they range up to $9{,}999{,}999$. The device page is $4$\,KiB. The measured workload is:
\begin{itemize}
  \item[(i)] $70\%$ point lookups by premises ID from field crews;
  \item[(ii)] $30\%$ monthly exception reports listing every premises in a pincode-and-street range, in ID order --- a typical report covers $2{,}500$ consecutive IDs;
  \item[(iii)] $3\%$ annual growth, with new IDs arriving in increasing order;
  \item[(iv)] one end-of-year pass that reads every record to compute totals.
\end{itemize}
\begin{enumerate}
  \item[(a)] Compute the records per block and the number of data blocks. Then design an \emph{indexed-sequential} organization: give the block layout, a sparse index with one entry per block (highest key plus block address, $16$ bytes per entry), the number of index entries, the entries per index block, the number of index blocks, and the number of block reads for one point lookup. Draw the three levels.
  \item[(b)] Using your layout, give the block reads for (i) one point lookup, (ii) one $2{,}500$-premises report, and (iii) the end-of-year pass. Then compare: the same point lookup on a sorted sequential file (binary search over blocks), and the same range report on a hashed file. Give both numbers.
  \item[(c)] Justify the indexed-sequential choice against all four alternatives (pile, sorted sequential, direct, hashed) on \emph{this} workload, using the numbers from (b) and the sparse key range. State the insert/overflow plan for the $3\%$ growth, when reorganization is triggered, and its cost in block transfers.
\end{enumerate}
Why this matters: the mixed point-plus-range workload is the normal case in practice, and it is the one workload where the fastest point-lookup structure is the wrong answer --- this question is the rehearsal for choosing a real index.

\bibliography{../../Bibliography_base}
\bibliographystyle{plain}
\end{document}
